\input{hand-preamble.tex}
\begin{document}\shapka
\lk{4}{Дерево формул. Дерево $N$-вывода}
\opr \term{Деревом формул} высоты $h$ с мн-вом листьев $\Gamma$ и корнем $F$ наз.:\par
1. Всякая формула $F$ есть дерево формул высоты $h=0$ с корнем $F$.\par
2. Если $D_1,\dots,D_n$ --- деревья ф-л, $h_1,\dots,h_n$ --- их высоты, то\par
\hspace*{20mm}$\dfrac{D_1\ \dots\ D_n}{F}$ --- дерево формул.\par
3. Других деревьев формул нет.\par
\red{Зам.} Отношение $N$-выводимости рефлексивно, монотонно и транзитивно.\par
\term{Утв.} Если $\Gamma\vdash_N A$ и $\Gamma\vdash_N B$, то $\Gamma\vdash_N A\&B$.\par
\dvo Поскольку всякое $N_{int}$-правило явл-ся $N_{cl}$-правилом, то всякий $N_{int}$-вывод явл-ся $N_{cl}$-выводом. \ctd\par
Пример. Построим дерево $N$-вывода формулы $\neg(A\vee B)\supset(\neg A\&\neg B)$:\par
\begin{prooftree}\AxiomC{$[A]^2$}\lab{$\vee$в}\UnaryInfC{$A\vee B$}\AxiomC{$[\neg(A\vee B)]^1$}\lab{$\neg$в(2)}\BinaryInfC{$\neg A$}\end{prooftree}
Семантическое следование (сс). Будем говорить, что из мн-ва формул $\Gamma$ семантически следует формула $F$ (и писать $\Gamma\models F$), если на любом ист. наборе, на котором каждая формула из $\Gamma$ принимает знач. И, формула $F$ тоже принимает знач. И.\par
\framebox{МЛ с.37}\par
\end{document}
